\section{Active globalization with exact marking}\label{ap:active}

We now return a score through an active layer $j\leq s$.  The local
orbit includes the higher active determining slots, so its replacement must
retain the formal identity $\pi^0g=P_Y$ from
Lemma~\ref{ap:ascent:representation}.  This identity will restore the
old integer arguments exactly in Section~\ref{ap:extraction}.

The passive proof could restrict the scoring slice whenever it froze a
slow factor or a rational factor.  Here the auxiliary bounds may depend
on $Q_j$.  Such a restriction could reduce the score by a cold factor,
which the eventual relative width loss cannot afford.  We may select a
cold-small fraction of paths on which one construction works, but on
each selected path we must keep the original slice and its warm score.

We therefore retain short integer parameters in polynomial families
instead of fixing their values in the score.  We also retain the slow
and rational factors as polynomials in two sets of labels: normalized
real labels and integer labels.  On an actual path both describe the
same parameters.  In the forecasting functional $\mathcal K$, the two
parallel labels are separate integration variables, with independent
base measures before the incidence and output weights are imposed;
each inactive label and its normalization still come from one actual
integer.  The argument below
preserves this distinction through the degree induction.

\subsection{The inductive statement}

Write $t=(v,z)$ for the current parameter variables.  The entries of
$z$ are distinguished short variables: they are actual integer
variables among the inactive parameters $t_I$, and their lengths have
pre-cost.  The list can initially be empty.  The other variables $v$
are currently long.  Further inactive variables may become short
during the argument.  A parallel variable never becomes short,
because its length $L$ is chosen after all the pre-thresholds below.
We retain the labels
\[
 \theta=(t_i/T_i)_i,\qquad \widetilde t=t
\]
on an actual path.  In the forecasting measure $\mathcal K$, these
mean separate normalized real and residue labels for the parallel
variables, and the actual integer inactive variables with their
normalizations.  All references to residue representatives below
concern only the parallel integer labels.

The induction changes the functions tested on the group as well as the
orbit.  We need two different finite approximations for these functions.
At one site, a point-value mesh approximates the nilmanifold function
itself.  Across sites, a field mesh approximates one path-indexed field
uniformly over the parent box.  The latter lets us choose a controlled
list of fields first; binning the values of each chosen field then
gives fixed site masks to which the sampler applies.  Neither assertion
requires regularity in the site variable, or a mesh of orbit
coefficients.  We state the precise conditions using the quotient
metrics and complexity conventions of Definition~\ref{prelim:niltest}.

\begin{definition}[Controlled external fields]\label{ap:active:fields}
Let $G^0/\Gamma^0$ have controlled rational filtered data.  A pair of
fields
\[
 F^{\pm}_{\psi,u}(\theta,\widetilde t)
       :G^0/\Gamma^0\longrightarrow[0,1]
\]
is controlled if it has the following properties at every requested
pre-accuracy.
\begin{enumerate}
\item Its Lipschitz constants on the quotient, and in the normalized
real labels with all other arguments fixed, have pre-cost.  In the
long integer labels it is periodic with a common pre-bounded modulus
for each path.  The moduli belong to a pre-bounded list and divide the
eventual modulus $\mathfrak M$.
\item The point-value family over sites, paths, and admissible labels
has a supremum-norm mesh of pre-logarithmic size.
\item For fixed short labels and fixed mesh and residue arguments in
the remaining labels, the family indexed by $\psi$ has a mesh in
supremum norm over the sites, again of pre-logarithmic size.
\end{enumerate}
The bounds may charge the accuracy, the auxiliary dimensions, and
the short lengths.  The assertion concerns only the chart and label
ranges in use.  Integer labels can be represented by arbitrary
integers.  No mesh of polynomial orbit coefficients is required.
\end{definition}

The initial fields of Lemma~\ref{ap:ascent:representation}
satisfy this condition.  They include the passive site multipliers
and, on the positive side, the input $f_0$, but exclude the current
flag, slice, and target.

At a positive degree $r$, the degree reduction is governed by
Lemma~\ref{lifting:refiltration}.  For a strong projection
$H\to H_F$ that is injective on its top layer, an element $h$ is
uniquely determined by its lower-degree image $hH_r$ and its marked
image in $H_F$.  We will arrange this situation for the middle factor
of the local orbit.  The fixed marked polynomial then supplies the
second coordinate, leaving a problem of degree $r-1$.

The original projection need not have this injectivity.  Our first
step uses the top frequencies visible to the payoffs to construct a
rational graded algebra whose top projects injectively.  Our second
step converts that symbol information into exact factorizations over
one common marked input.  The last step passes to the lower quotient
and restores the original two payoffs by averaging.  Throughout this
construction, restrictions used to discover identities change the
selected paths, not the slice in the objective score.

\begin{theorem}[Active replacement]\label{ap:active:replacement}
Fix the prepared parent and a feasible center at a layer $j\leq s$.
Assume the modified sampler, scalar comparison, and major-decomposition
contracts established above.  Let $G^0$ be a controlled filtered
group of degree $r$, with a strong rational projection
$\pi^0:G^0\to Y$.  Suppose a pre-inverse fraction of the available
paths are modular-good in the sense of
Equation~\eqref{ap:scalar:modular-goodness}, and each such path has a
polynomial $g_\psi(v,z)$ whose logarithmic coordinates
have bounded total degree, adapted in $v$ at each fixed $z$, and
satisfy the formal identity
\begin{equation}
 \pi^0g_\psi(v,z)=P_Y(X_\psi(v,z);z),
 \qquad X_\psi=(\psi,\beta\circ\psi).
 \label{ap:active:initial-marking}
\end{equation}
Here $P_Y(X;z)$ is fixed on those paths, weighted filtered in $X$,
and polynomial in $z$, with bounded total degree.  There is no bound
on the coefficients of either polynomial.

Let the two external fields satisfy
Definition~\ref{ap:active:fields}.  Suppose their signed score, with
the original slice mask $1_{\mathcal S}$ and full parameter mass, is
at least $\delta>0$:
\begin{equation}
 \E_t 1_{\mathcal S}(t)
 \left[
  1_{\mathrm{no\ bad}}(t)
  F^-_{\psi,\psi(t)}(\theta,\widetilde t)
       (g_\psi(t)\Gamma^0)
  -\lambda F^+_{\psi,\psi(t)}(\theta,\widetilde t)
       (g_\psi(t)\Gamma^0)
 \right]\geq\delta,
 \label{ap:active:input-score}
\end{equation}
where $0<\lambda\leq1$.  Include the original slice cost, the short-label
range bounds, and $\log(1/\delta)$ in the prescribed preliminary data.
All path probabilities refer to the original ensemble, with discarded
paths retained as zero mass; subsequent selections are subsets of the
modular-good productive paths just specified.  For any prescribed
$0<\eta<\delta/2$, include $\log(1/\eta)$ in that request.  Then one
retained modular-good path admits a stochastic polynomial
\begin{equation}
 g^\sharp(X,\theta,\widetilde t)\in G^0,
 \qquad \pi^0g^\sharp(X,\theta,\widetilde t)=P_Y(X;z)
 \label{ap:active:output-marking}
\end{equation}
whose $\mathcal K$-score, with the same slice and flag masks and the
expectation of the two payoffs, is at least $\delta-\eta$.

A stochastic polynomial means a polynomial formula with random
constant parameters, sampled independently of its arguments.  Its
total logarithmic-coordinate degree is bounded in terms of the
input degree bounds only; its coefficients need not be bounded.  A
filtered-degree bound in the original parameters is not asserted.
On the joint chart and exact-inactive locus
\begin{equation}
                    z_I(X)=z_I(t_I),
 \label{ap:active:exact-locus}
\end{equation}
each expected payoff separately is independent of the representatives
chosen for the parallel residue labels, with all other data fixed.
This is a pointwise statement for every admissible label choice.

The retained path fraction and all auxiliary logarithmic costs may
have cold pre-cost.  The number of successive polynomial budget
increases and all degree increases depend only on the fixed input
degrees.  Thus, when $\delta$, the original slice, and $\eta$ have
warm bounds, the resulting score and slice still have those bounds.
\end{theorem}

\begin{proof}
We induct on $r$.  Every loss of score below is a prescribed additive
part of $\eta$.  There are boundedly many such losses at each of the
boundedly many degree drops.  Productive path fractions, unlike the
score in~\eqref{ap:active:input-score}, may decrease by pre factors.
We always budget their inverse logarithms before the next modeling
or separation request.  Every retained subset consists of modular-good
paths, so the eventual scalar comparison receives its stated goodness
hypothesis even when the retained fraction has cold preliminary cost.

If $r=0$, the group is trivial.  At fixed short labels, subdivide
normalized long-label ranges and integer residue classes according
to the controlled field meshes.  On every piece the same bounded
site function approximates the payoff both on an actual path and
on the corresponding forecasting support.  Axes shorter than the
mesh and residue cutoffs are enumerated exactly; all other pieces
have pre-relative sides.  Parallel sides exceed these cutoffs, so
endpoint rounding has negligible mass.  Apply
Proposition~\ref{ap:scalar:comparison}, including its flag version,
to the fixed menu of site functions.  A union bound over menu
entries and testing slices gives the claimed comparison on retained
paths.  We do not pigeonhole an arbitrary list of mesh choices
across all sites.  The trivial polynomial is the output, and the
representative assertion follows from periodicity of the fields.

Assume henceforth $r>0$.

\subsubsection*{1. Constructing an injective top projection}

\paragraph{Keeping short parameters in a polynomial family.}

Use compatible rational splittings of the filtrations of $G^0$ and
$Y$.  Form the Lie group $\widehat G^0$ of logarithmic polynomial
functions in dummy short variables $z'$: in a split original
layer~$i$, impose a degree cap $Ci$ in $z'$.  The filtration is by
pointwise membership in the original layers.  Choose the fixed
integer $C$ large enough for the maps below.  A bracket of layer-$i$
and layer-$h$ coordinates has dummy degree at most $C(i+h)$, so
the caps are closed under BCH.  Thus $\widehat G^0$ is a rational
nilpotent Lie group of polynomial dimension.  Compatible sections
give a strong projection to the analogous group $\widehat Y$.
Evaluation at any specified $z'$ is strong before refiltration.

Choose one controlled deep lattice $\widehat\Gamma^0$ whose
evaluation at every integral $z'$ lies in $\Gamma^0$.  Indeed a deep
rational integral grid of monomial coefficients has this property,
and the bounded-step rational group operations preserve it after
one further fixed divisibility.  This is a simultaneous lattice
choice, not a separate choice at each short label.

The families $z'\mapsto g_\psi(v,z')$ and
$z'\mapsto P_Y(X;z')$ fit the caps.  The common-site correction used
below also fits: it uses only the given maps, whose site components
have degrees at most $1$ and $j$, a pointwise filtered logarithmic
section, and bounded-step BCH.  Consequently $C$ can be fixed from
the input degrees before partner selection.

There are now two site inputs to distinguish.  At an actual short
label $z$, the local family has projected value
\[
 z'\longmapsto P_Y(X_\psi(v,z');z'),
\]
whereas reconstruction at the actual site $X_\psi(v,z)$ will require
\[
 z'\longmapsto P_Y(X_\psi(v,z);z').
\]
These families agree at $z'=z$.  Their pure-long leading symbols
also agree, because the leading site input is independent of the short
labels.  Their full polynomial maps need not agree.  We may therefore
use the uncorrected family to find the top algebra now; after matching
the fixed marked polynomial, we will correct the full family at the
actual site.

\paragraph{Averaging over the top kernel.}
Apply Lemma~\ref{prelim:vertical} in the pointwise top torus $G^0_r$
to both external functions.  Use a fixed convolution and Fourier
expansion with uniform error smaller than the assigned tolerance.
Taking a retained vertical component preserves controlled binning,
with pre amplification of approximation errors.  At each actual
short label $z$, call a truncated mode significant when the
absolute value of its conditional tested average over the remaining
variables, including that term's masks, exceeds the tolerance
divided by a sufficiently large constant times the number of modes.
The sum of discarded averages, integrated over all short labels
with their original masses, is then negligible.

Let $L_0$ be the connected joint annihilator, in
$\widehat G^0_r\cap\ker\widehat\pi^0$, of the significant pointwise
frequency rows composed with evaluation at their respective $z$.
These rows are integral on the family lattice.  Pigeonhole a basis
of them, together with its frequencies, short labels, and payoff
choices.  There are polynomially many basis rows and a pre number
of possible basis lists.  In particular, we do not enumerate all
subsets of the truncated modes.  The rational central subgroup
$L_0$ is now fixed on a pre-inverse fraction of paths.

Average both payoffs by left translation over
$L_0/(L_0\cap\widehat\Gamma^0)$, evaluating a family translate at
the actual short label.  A character is unchanged if its row
vanishes on $L_0$ and is killed otherwise.  Hence every significant
mode survives, and the full score on the original slice changes by
only the discarded-average and Fourier-tail errors.  Both averaged
functions remain in $[0,1]$.

\paragraph{Partners and the joint fast algebra.}

For each pivot row we obtain a step drop paired with a fixed
ambient partner.  We give the selection details because these
auxiliary restrictions must not be imposed on the objective score.

If the pivot comes from the positive flagged term, first truncate
the inclusion-exclusion expansion of the no-bad flag as in the
proof of Proposition~\ref{ap:scalar:comparison}, using
Lemma~\ref{ap:scalar:flags}.  The error may charge the short
label mass, mode count, and current productive probability.  Split
each retained term into its pre residue subpieces, so its coefficient
is constant on each piece for the chosen path.

At the pivot's short label, subdivide the remaining normalized labels
and integer classes for controlled binning.  First approximate the
site field by its path-field mesh.  Then bin its point values using
the point-value mesh.  The triangle inequality supplies a biased
pairing of a fixed vertical function with an arbitrary bounded
masked site input from a fixed pre menu.  Pieces too small to
support the bias can be discarded for this selection.  All slices
used for detection therefore have pre-relative sides, including
the specialization of the short variables.  Their orbit is still
filtered of degree $r$ in the remaining parameters.

Use the sampler's pathwise testing-seminorm approximation at
degree $r$ on all these masked inputs.  Markov's inequality and a
union bound pay for the menu and the current productive probability.
One model term gives correlation with a fixed ambient degree-$r$
partner $Q_i(u)$, its niltest function, and a site twist.  Truncate
the partner's top Fourier expansion, freeze the twist on a finer
slice, and select a biased pair of vertical functions.  Freezing is
permitted because normalized spatial and lift coordinates have
controlled path coefficients and the lift residues can be fixed.
Freeze further axes that are too short for
Lemma~\ref{prelim:step-drop}.  That Lemma gives a splitting of the
pair's pure-long symbol into a rational fast algebra annihilated by
the joint top frequency.  Translation, restriction, and dilation
refer this homogeneous identity back to the original parameters
with pre-cost.

Perform these selections simultaneously for the polynomially many
pivots, exactly as in the partner-selection part of
Proposition~\ref{lifting:candidate-globalization}.  Pigeonhole
partners, top rows, and the controlled-height rational fast algebras.
The input test complexities in the modeling request were fixed
before partners were selected.  Symbol slices for distinct pivots
need not intersect: we use only their common remaining pure-long
variables.  At each original layer the total degree in the previous
long variables is bounded by that layer, so specializing newly
short variables or changing their origins cannot change these
pure-long top terms.  Jointly raise any interpolation and separation
thresholds after the selections.  No new slice is imposed on
\eqref{ap:active:input-score}.  If there are no pivots, use no
partners and the full fast algebra.

The pure-long leading site input
$X_{\mathrm{top}}=(\psi_{[1]},\beta_{[j]})$ is independent of all
short labels.  Therefore evaluation and coordinate restriction of
the joint family symbol give exactly the pair symbols just obtained.
Pull each pair algebra back under those strong maps, lift its outer
factors, and join the splittings using
Lemma~\ref{prelim:symbol-splitting}.  After quotienting by $L_0$, set
\begin{equation}
 \begin{aligned}
 G&=(\widehat G^0/L_0)\times\prod_iG_{Q_i},
 &F&=\widehat Y\times\prod_iG_{Q_i},\\
 \pi&:G\longrightarrow F,
 &P_F(X)&=\bigl((z'\mapsto P_Y(X;z')),(Q_i(u))_i\bigr).
 \end{aligned}
 \label{ap:active:joint-groups}
\end{equation}
Let $\mathfrak v\subset\operatorname{gr}G$ be the joined fast
algebra and $\mathfrak v_F=\operatorname{gr}\pi(\mathfrak v)$.
The constraints from the partners now give the required injectivity:
\[
               \mathfrak v_r\longrightarrow(\mathfrak v_F)_r
               \quad\hbox{is injective}.
\]
Indeed a top vector in its kernel has zero partner coordinates;
its family component in the marking kernel is annihilated by every
pivot and thus lies in $L_0$ before taking the quotient.  The marked
projection of the joint symbol is the weighted symbol of $P_F$
evaluated on $X_{\mathrm{top}}$.  Thus the top-layer ambiguity has been removed.  The remaining work
before degree reduction is to factor the full maps with common outer
projections.  The marked identity established so far concerns only
symbols: the two full family projections displayed above still have
to be reconciled.

\subsubsection*{2. Exact factors over one common site}

\paragraph{Matching the fixed marked symbol.}

Apply the marked-symbol argument of
Proposition~\ref{lifting:candidate-globalization} in
\eqref{ap:active:joint-groups}.  It gives slow and rational outer
symbols $\mathcal E,\mathcal R$ for $P_F$ on a controlled-height
tag space $K\supset W$ in the $\beta$ variable, with residual in
$\exp\mathfrak v_F$.  For completeness, at each grade raise the
long threshold so that the leading $\beta$-evaluation lies in the
currently imposed space.  Its projection off that space has
bounded-denominator rational coefficients, because $\beta$ is
integral, and is fully slow, because $C-c\in W$.  Separation makes
the pure-long coefficients vanish.  Compare the projected local
splitting with the already matched grades, reset their outer
factors through $\mathfrak v_F$, and pass to the quotient at the
next grade.  The remaining fixed scalar coordinates evaluate to
slow plus rational polynomials.  Proposition~\ref{lifting:candidate-major},
with the cold detection and chart identities verified above,
supplies their ambient decompositions.  Intersect the resulting
spaces and lift the corrections by rational sections.  The slow
coordinates are bounded in $u/H,\beta-A(u)$, where $A=C_{[j]}$;
the rational ones have controlled common denominator in raw
arguments.  There are only $r$ grades.  Polynomial families in
$\widehat Y$ are coordinates in a fixed controlled rational group,
so this is the same finite grade induction, not an induction over
their individual coefficients.

Raise the long threshold after these outputs, and let
$z_{\mathrm{all}}=(z,z_{\mathrm{new}})$ be all short coordinates,
with remaining variables $v_L$.  The whole off-$K$ polynomial has
the form
\begin{equation}
 \operatorname{proj}_{K^\perp}\beta(\psi(v_L,z_{\mathrm{all}}))
                      =o(z_{\mathrm{all}}).
 \label{ap:active:off-space}
\end{equation}
Its coefficients and common denominator have pre bounds.  Indeed
each coefficient involving a long variable is rational with pre
denominator and is smaller than its reciprocal denominator by the
fully slow path bound, hence is zero.  Pigeonhole the remaining
coefficient list.  Define
\begin{equation}
 \begin{split}
 p(X,z_{\mathrm{all}})
   &=(u,\operatorname{proj}_K\beta+o(z_{\mathrm{all}})),\\
 P_F^*(X,z_{\mathrm{all}})&=P_F(p(X,z_{\mathrm{all}})).
 \end{split}
 \label{ap:active:projected-input}
\end{equation}
On path data $p(X_\psi,z_{\mathrm{all}})=X_\psi$ polynomially.
The same equality holds on the forecasting locus
\eqref{ap:active:exact-locus}.  To see this, choose pre-height
integer annihilating rows for $K$.  They vanish on $W$ and pair
integrally with the projected lattice basis vectors.  Their values
on active sheet vectors are zero: those vectors have size
$O(1/K_i)$, with $k_i\gg(AL)^j$, and the threshold exceeds the row
height.  The row values therefore depend only on the inactive sheet
coordinates, which match exactly.  The short labels are actual
entries of the same $t_I$.

Refilter in the inverse images $H,H_F$ of
$\mathfrak v,\mathfrak v_F$, using
Lemma~\ref{lifting:refiltration}.  Lift the matched symbols to full
factors
\begin{equation}
                 P_F^*=E_F P_H R_F,\qquad P_H\in H_F.
 \label{ap:active:marked-full-factors}
\end{equation}
In the slow factor use $u/H$ and
$\operatorname{proj}_K\beta-(C-c)(u)$; in the rational factor use
$u,\operatorname{proj}_K\beta$.  Their middle symbol belongs to
$\mathfrak v_F$, so positive-weight coefficients are adapted to
the new filtration.  In the vector quotient $F/H_{F,1}$, only the
constant in $X$ remains.  It is polynomial in $z_{\mathrm{all}}$.
Split its coefficients, in rational coordinates, into bounded and
integral parts and insert lifts on the inside of the outer factors.
This is a coefficientwise split, not taking fractional parts of a
polynomial at varying labels.  It preserves its degree.  We obtain
$E_F$ fully bounded with pre coefficients in the indicated centered
and normalized arguments and short labels, $R_F$ with pre common
denominator in raw integer arguments and short labels, and $P_H$
weighted filtered in $X$.  These factors are fixed across the
retained paths.

\paragraph{Correcting the full family marking.}

We must correct the full family marking before making local factors.
Put
\[
 X_a=X_\psi(v_L,z,z_{\mathrm{new}}),\qquad
 X_{z'}=X_\psi(v_L,z',z_{\mathrm{new}}).
\]
Let $\rho^0$ be a filtered logarithmic section of the pointwise
projection $\pi^0$.  Before quotienting the family, replace its
local map by
\begin{equation}
 z'\longmapsto
 \rho^0\!\left(P_Y(X_a;z')P_Y(X_{z'};z')^{-1}\right)
                  g_\psi(v_L,z',z_{\mathrm{new}}),
 \label{ap:active:common-site-correction}
\end{equation}
and use the partners $Q_i(\psi(v_L,z,z_{\mathrm{new}}))$.
Call the resulting map into $G$ $\mathbf g_\psi$.
Its family value at $z'=z$ is unchanged.  The inserted ratio has
identity pure-long symbol, since $X_a$ and $X_{z'}$ have the same
pure-long leading input.  Thus the former joint symbol splitting
is retained as a whole family, independently of the values of
$z_{\mathrm{all}}$.  But its full projection is now
$P_F^*(X_a,z_{\mathrm{all}})$, at one common actual input.  Explicitly,
the projected family product is
\[
 [P_Y(X_a;z')P_Y(X_{z'};z')^{-1}]P_Y(X_{z'};z')
                         =P_Y(X_a;z').
\]
No commutativity is needed.  This correction is why we do not need
evaluation to commute with refiltration.

\paragraph{Lifting the prescribed outer factors.}

We construct
\begin{equation}
 \mathbf g_\psi=E_\psi h_\psi R_\psi,
 \qquad \pi E_\psi=E_F,\quad
 \pi R_\psi=R_F,\quad \pi h_\psi=P_H,
 \label{ap:active:local-full-factors}
\end{equation}
where projected factors use $X_a,z_{\mathrm{all}}$ and $h_\psi$ is
adapted to $H$ in $v_L$.  More importantly, $E_\psi,R_\psi$ have
ambient-with-label versions with these exact formal projections.
They are products of fixed logarithmic-section lifts of $E_F,R_F$
and kernel discrepancies.  The slow discrepancies have bounded
polynomial coefficients in normalized long labels $\theta_L$ and
short labels; the rational discrepancies have controlled common
denominator in integer labels.

Here are the exact corrections.  Write the local symbol splitting
as $ePd$, with $P\in\exp\mathfrak v$.  Let $e_F,d_F$ be the
symbols of the prescribed projected outer factors.  By compatible
splitting, on sufficiently long axes,
\[
       (\pi e)^{-1}e_F\in\exp\mathfrak v_F,
       \qquad d_F(\pi d)^{-1}\in\exp\mathfrak v_F.
\]
The first is slow and the second rational.  Lift them through a
controlled graded section inside $\mathfrak v$ and reset the local
splitting.  Its outer symbols now project exactly to $e_F,d_F$.
Fix a filtered logarithmic section $\rho$ of $\pi$, and denote the
reset outer symbols by $e',d'$.  The two ratios
\[
 e'\bigl(\operatorname{Sym}\rho(E_F)\bigr)^{-1},
 \qquad
 \bigl(\operatorname{Sym}\rho(R_F)\bigr)^{-1}d'
\]
are kernel symbols.  Strongness gives actual
kernel lifts: if a representative in the $i$th layer projects into
the $(i+1)$st, subtract a section lift in that deeper layer.
Coefficientwise lifts yield kernel polynomial maps on the left of
$\rho(E_F)$ and on the right of $\rho(R_F)$, respectively, with
the claimed bounds.  Their residual has all positive-long-degree
coefficients adapted to the layers of $H$.

For its constant in $v_L$, use the vector quotient $G/H_1$.
Its image belongs coefficientwise to the image of $\ker\pi$,
because its projection is supplied by $P_H\in H_F$ and
$\pi(H)=H_F$.  Split that short-label polynomial coefficientwise
into bounded and integral parts, lift within the actual kernel,
and insert the factors on the inside of the two outer factors.
This puts the remaining constant in $H_1$, preserves all exact
projections, and proves~\eqref{ap:active:local-full-factors}.

We choose the ambient versions after this constant adjustment in
the definite order
\begin{equation}
 \begin{aligned}
 E_\psi&=K_{E,\psi}(\theta_L,z_{\mathrm{all}})
                   \rho(E_F(X,z_{\mathrm{all}})),\\
 R_\psi&=\rho(R_F(X,z_{\mathrm{all}}))
                   K_{R,\psi}(\widetilde t_L,z_{\mathrm{all}}).
 \end{aligned}
 \label{ap:active:ordered-outer-factors}
\end{equation}
The kernel polynomials here have no additional site argument.  This
order will let us compare path-dependent fields while the fixed site
factors cancel: at a fixed site and labels,
\[
 R_\psi^{-1}R_{\psi'}=K_{R,\psi}^{-1}K_{R,\psi'}.
\]
Thus unrestricted coefficients of $R_F$ will not enter a mesh of the
rational discrepancies.  To preserve this order when absorbing an inside constant, first
conjugate it past the common factor \emph{on the path}, and express
the conjugation using the pullback
$\rho(E_F(X_\psi(t),z_{\mathrm{all}}))$ or
$\rho(R_F(X_\psi(t),z_{\mathrm{all}}))$ as a polynomial in the
labels.  The first pullback is fully slow with controlled normalized
coefficients.  The second has controlled common denominator in
integer labels, because $X_\psi$ has integral polynomial
coefficients.  Extend these kernel label formulas off the path,
while retaining the displayed common site factors.  They still
agree with the corrected local factors on paths and have the
exact formal projections everywhere.  This procedure bounds
denominators independently of the unrestricted integral
coefficients of $X_\psi$.

The threshold order here is important.  Any final increases needed
for separation compare the pure-long family symbols with the
already constructed matched symbols.  Those symbols have no cost
from lengths of axes made short by this very increase.  Likewise
off-space vanishing uses full coefficient slowness of $y$.
Subsequent polynomial estimates may charge the resulting short
lengths.  Thus there is no threshold iteration once per axis.

At this point the full factorization~\eqref{ap:active:local-full-factors}
has the exact projections required for reconstruction, and its middle
map is adapted to the refiltered group $H$.  The first stage proved
injectivity at the top.  We can now remove that top layer while
retaining the two payoffs as functions on the lower quotient.

\subsubsection*{3. Degree reduction and restoration of the score}

\paragraph{Reconstructing the payoff on the lower quotient.}

Read the $L_0$-averaged original fields on $G$, evaluating a lift of
the family quotient at its actual diagonal short label and ignoring
partners.  Denote them $\overline F^{\pm}_{\psi,u}$, with label
arguments suppressed.  They are well defined because $L_0$ is
central and the family lattice evaluates into the pointwise lattice
at integral labels.  Their Lipschitz bounds have pre-cost, using
bounded-label evaluation and controlled quotient sections.  This
averaging is nonexpansive in supremum norm at fixed labels.
By~\eqref{ap:active:common-site-correction}, their local evaluations
on $\mathbf g_\psi$ are exactly the retained averaged payoffs.

Consider, at each site and label, the function on $H$
\[
       h\longmapsto
       \overline F^{\pm}_{\psi,u}(E_\psi hR_\psi),
\]
using the ambient-with-label factors.  It descends on a controlled
deep source lattice.  Top injectivity and
Lemma~\ref{lifting:refiltration} give a rational filtered isomorphism
\begin{equation}
 H\simeq (H/H_r)\mathop{\times}_{H_F/H_{F,r}}H_F.
 \label{ap:active:fiber}
\end{equation}
Choose sufficiently deep lattices in the two target groups so
their inverse image lies in the source lattice.  In logarithmic
coordinates this follows by applying the rational linear inverse
of the joint embedding to a sufficiently divisible product grid.
The reconstruction argument of
Proposition~\ref{lifting:candidate-globalization} therefore gives
a common controlled Lipschitz constant on the joint image.  Extend
each function by the same clipped infimum extension, and specialize
the second coordinate at $P_H(X,z_{\mathrm{all}})$.
The extension agrees exactly on compatible pairs.

For clarity, a common source grid permits all right translates
$R_\psi$: their logarithmic coordinates have one controlled common
denominator after pigeonholing denominator values, and the adjoint
formulas have bounded denominators independently of numerators.
Left translation has controlled Lipschitz cost because $E_\psi$
is fully bounded.  These are the hypotheses required by the
lattice reconstruction, not bounds on the middle orbit.

\paragraph{Preserving the two mesh conditions.}
The resulting fields on $H/H_r$ satisfy
Definition~\ref{ap:active:fields}.  For point-value meshes, the
parameters are the old function value, bounded $E_\psi$, the right
lattice coset of $R_\psi$, bounded-label evaluation, and the compact
specialization value of $P_H$.  The clipped infimum extension is
nonexpansive in its input function and Lipschitz in specialization.
Thus these are finite-dimensional parameter meshes; no net for all
Lipschitz functions is used.

For field meshes in the path index, the marked functions
$P_H,E_F,R_F$ are fixed as functions of the site and short labels.
Slow kernel factors have controlled coefficient meshes.  Rational
kernel factors have a pre number of right-coset fields despite
unrestricted numerators.  Indeed prescribe their polynomial
coefficients modulo a sufficiently divisible pre grid.  In the
BCH expression for the difference, every noncancelled term contains
a coefficient difference, and all other denominators are controlled;
at integer labels the difference belongs to the prescribed deep
lattice.  This fixes the right coset uniformly over sites and
integer labels.  The same calculation supplies a pre period in
long integer labels.  Taking a common selected period preserves
divisibility by the eventual $\mathfrak M$.
The cancellation following~\eqref{ap:active:ordered-outer-factors}
shows that this comparison involves only the kernel label factors.
The common site factor $\rho(R_F)$ introduces no conjugation or
coefficient bound.

Continuous long-label variation enters only the old fields and
the controlled slow factors.  It does not enter the possibly
unbounded fixed site polynomial $P_H$, which uses the actual site
and exact short labels.  Additional short axes contribute only
pre many label choices to the meshes.  These observations verify
all parts of the refreshed field contract without meshing any
unrestricted orbit coefficient.

\paragraph{The lower-degree return and stochastic restoration.}

We have a lower-degree problem on
$H/H_r\to H_F/H_{F,r}$ with actual local orbit
$h_\psi\bmod H_r$, exact marking
$P_H(X,z_{\mathrm{all}})\bmod H_{F,r}$, and the refreshed fields.
Its score uses exactly the original slice and flag.  Apply the
induction hypothesis with the remaining error tolerance.
Its stochastic output lifts jointly with the exact $P_H$ to a
polynomial $h^\sharp\in H$ by~\eqref{ap:active:fiber}.  Multiply
$E_\psi h^\sharp R_\psi$ on the left by
\begin{equation}
       \rho\!\left(P_F(X)P_F^*(X,z_{\mathrm{all}})^{-1}\right).
 \label{ap:active:restore-marking}
\end{equation}
The projection is now exactly $P_F(X)$ formally.  The multiplier
is the identity on~\eqref{ap:active:exact-locus}, by
\eqref{ap:active:projected-input}.  Project to the family quotient,
lift by a logarithmic section, multiply by an independent uniform
constant in $L_0/(L_0\cap\widehat\Gamma^0)$, using bounded Borel
representatives, and evaluate at $z'=z$.  This defines
$g^\sharp$ and proves its formal marking
\eqref{ap:active:output-marking}.

Fix any admissible site and labels on the exact-inactive locus.
The reconstructed pair is compatible, so the extended lower-degree
function agrees exactly with the translated averaged original
function.  The correction~\eqref{ap:active:restore-marking} is the
identity there.  Averaging the newest family translate realizes
exactly the earlier $L_0$-average after diagonal evaluation.  Haar
invariance removes any dependence on the selected family lift,
and evaluation of its lattice lies in the original pointwise
lattice.  Consequently, for each label choice, the expectation of
each restored payoff equals the corresponding lower-step payoff.
Taking the remaining expectations gives representative independence
by induction, separately for the two signs.  It also preserves
the $\mathcal K$-score.  This is an identity of expectations;
an individual restored polynomial need not be periodic in the
parallel integer labels.

All random parameters just introduced are constants independent of
the arguments.  Evaluation of a fixed random polynomial family at
the variable short label remains polynomial.  Group multiplication,
logarithmic sections, the common-site correction, coefficient lifts,
and family evaluation are bounded-degree operations.  They are
nested only a bounded number of times in this degree induction.
Taking simultaneous partners, solving all component equations,
or intersecting many subspaces does not iterate a degree increase
once per component.  Hence total degree depends only on the input
degree bounds, while dimensions and pre logarithmic costs have
boundedly many successive polynomial increases.

Finally, forecast the flag truncation, modeling, symbol, and period
bounds at each of those pre stages before choosing $Z$.  Choose
$Z$ to include the required individual periods, and only then
choose perturbation, range, rank, and side cutoffs.  The scalar
comparison uses individual moduli and witness products, not a
Fourier expansion at the full possibly enormous $\mathfrak M$.
All losses to the objective have been the prescribed additive
parts of $\eta$; short-label selection, partner discovery, and
symbol restrictions reduced only productive path fractions.
This proves the theorem.
\end{proof}

\begin{remark}[What the macro argument may use]\label{ap:active:macro-interface}
In the initial instance of Theorem~\ref{ap:active:replacement}, the
fields have no direct parameter-label dependence: the labels occur
only in the replacement polynomial.  All auxiliary label and period
subdivisions used inside its proof have disappeared from the score's
slice mask.  The macro argument may therefore change parallel
residue representatives pointwise, using the expected-payoff
identity, before fixing the stochastic constants by averaging.
It retains a warm full-mass score and the original warm slice, but
must not assume bounded coefficients or periodicity of an individual
replacement orbit.
\end{remark}
